\section{A triangular change of variables}
\label{sec:triangular}

An onto bi-Lipschitz Hilbert-space change of variables \(h\) will
produce a bi-Lipschitz equivalence between two weighted graph spaces.
The triangular map below has an explicit inverse obtained from
\(h^{-1}\). The identities involving \(h\) and \(h^{-1}\) control
the \(c_0\)-output, while direct Lipschitz estimates control the
weighted sum. Complete continuity of the linearization of
\(h-\pi_U\) enters only at the final step: the preceding obstruction
then shows that the two graph spaces are not linearly isomorphic.

Triangular changes of variables also appear in the quotient-section
construction of Aharoni and Lindenstrauss \cite[pp.~281--282]{AL78}
and in the free-space splitting of Godefroy and Kalton
\cite[Theorem~2.12]{GK03}. Here the argument uses the locally
verified identity \(qB=hQ_1\).

\begin{proposition}[Weighted graph criterion]
\label[proposition]{prop:criterion}
Let \(U\) be a separable real Hilbert space, \(V=\ell_2\), and
\(M=U\oplus_2V\), with projections \(\pi_U,\pi_V\).
Suppose \(h:M\to U\) is a bijection with \(h(0)=0\) and, for all
\(x,x'\in M\),
\[
 a\|x-x'\|\le\|h(x)-h(x')\|\le b\|x-x'\|,
 \qquad 0<a\le b<\infty.
\]
Let \(E=\F(M)\) and let \(q:E\to U\) be the linearization of
\(D=h-\pi_U\). Suppose \(q\) is completely continuous. Give
\(E\oplus V\) the sum norm, define
\[
 Q_1:E\oplus_1V\to M,\quad Q_1(e,v)=(qe,v),\qquad
 X=Z_{Q_1},\qquad Y=Z_q,
\]
and put \(L_q=\|q\|\), \(K_q=\max\{L_q,1\}\).
Then \(X,Y\) are separable real Banach spaces and there is a bijection
\(\Psi:X\to Y\) such that
\begin{equation}
\label{eq:criterion-bounds}
 \Lip(\Psi)\le\max\{1+K_q,b\},\qquad
 \Lip(\Psi^{-1})\le\max\{1+2L_q/a,1/a\}.
\end{equation}
Moreover, \(X\) contains \(c_0(V)\) isometrically and \(Y\) contains
no closed subspace linearly isomorphic to \(c_0(V)\).
\end{proposition}

\begin{proof}
The space \(E\) is separable because \(M\) is separable. The two
operators \(Q_1\) and \(q\) are bounded: in the sum norm,
\begin{equation}
\label{eq:q1-bound}
 \|Q_1(e,v)\|_M
 =\bigl(\|qe\|^2+\|v\|^2\bigr)^{1/2}
 \le K_q(\|e\|+\|v\|).
\end{equation}
Thus \Cref{lem:graph-banach} applies to \(X\) and \(Y\).

Define the zero-preserving map
\begin{equation}
\label{eq:B-definition}
 B:E\oplus_1V\to E,\qquad B(e,v)=e+\delta(qe,v).
\end{equation}
The defining identity \(q\delta(u,v)=h(u,v)-u\) gives
\begin{equation}
\label{eq:diagram}
 qB(e,v)=qe+h(qe,v)-qe=h(Q_1(e,v)).
\end{equation}
For \(z\in E\), the inverse of \(h\) is defined at \(qz\); set
\(x=h^{-1}(qz)\) and
\begin{equation}
\label{eq:B-inverse}
 C(z)=\bigl(z-\delta(x),\,\pi_Vx\bigr).
\end{equation}
Since
\[
 q(z-\delta(x))
 =qz-\bigl(h(x)-\pi_Ux\bigr)=\pi_Ux,
\]
we have \(Q_1C(z)=x\), and substitution in \eqref{eq:B-definition}
gives \(B(C(z))=z\). Conversely, if \(z=B(e,v)\), then
\eqref{eq:diagram} and injectivity of \(h\) give
\(h^{-1}(qz)=Q_1(e,v)=(qe,v)\). Formula~\eqref{eq:B-inverse}
therefore returns \((e,v)\). Thus \(C=B^{-1}\) and
\begin{equation}
\label{eq:inverse-diagram}
 Q_1B^{-1}=h^{-1}q.
\end{equation}
Only the inverse of \(h\) was used. In particular, this argument does
not require \(q\) to be onto.

We next keep track of the two norms. For \(w=(e,v)\) and
\(w'=(e',v')\), the isometry of \(\delta\) and
\eqref{eq:q1-bound} yield
\[
 \|B(w)-B(w')\|
 \le\|e-e'\|+\|Q_1(w-w')\|
 \le(1+K_q)\|w-w'\|_{E\oplus_1V}.
\]
For \(z,z'\in E\), put \(x=h^{-1}(qz)\),
\(x'=h^{-1}(qz')\). The lower bound for \(h\) implies
\[
 \|x-x'\|_M\le\frac{L_q}{a}\|z-z'\|.
\]
Using the \emph{sum} norm in \eqref{eq:B-inverse}, rather than a
Hilbert sum norm on \(E\oplus V\), gives
\begin{align*}
 \|B^{-1}(z)-B^{-1}(z')\|_{E\oplus_1V}
 &\le\|z-z'\|+\|\delta(x)-\delta(x')\|+\|\pi_V(x-x')\|\\
 &\le\|z-z'\|+2\|x-x'\|_M\\
 &\le(1+2L_q/a)\|z-z'\|.
\end{align*}
Neither estimate divides by \(L_q\); they remain valid as written
when \(L_q=0\). In the present setting \(V=\ell_2\ne0\), that case
cannot in fact occur under the hypotheses, since \(q=0\) would give
\(h(u,v)=u\), contradicting injectivity.

Define \(\Psi\) coordinatewise by
\begin{equation}
\label{eq:Psi-coordinatewise}
 \Psi(s)_i=B(s_i)\qquad(s\in X).
\end{equation}
Because \(B(0)=0\), the bound for \(B\) makes
\(\sum_i2^{-i}\|B(s_i)\|\) finite. Also \(Q_1s_i\to0\) in \(M\);
by \eqref{eq:diagram} and continuity of \(h\) at zero,
\(qB(s_i)=h(Q_1s_i)\to0\). Hence \(\Psi(s)\in Y\).
Likewise \(B^{-1}(0)=0\) controls the weighted sum for the
coordinatewise inverse, and \eqref{eq:inverse-diagram} shows that
if \(qz_i\to0\), then \(Q_1B^{-1}(z_i)=h^{-1}(qz_i)\to0\).
Thus the coordinatewise inverse maps \(Y\) into \(X\), and the
pointwise inverse identities make \(\Psi\) a bijection.

For clarity, the estimates apply to every difference of graph
sequences. For \(s,t\in X\),
\begin{align*}
 \sum_i2^{-i}\|B(s_i)-B(t_i)\|
 &\le(1+K_q)\sum_i2^{-i}\|s_i-t_i\|,\\
 \sup_i\|qB(s_i)-qB(t_i)\|
 &\le b\sup_i\|Q_1(s_i-t_i)\|.
\end{align*}
Taking the maximum proves the first inequality in
\eqref{eq:criterion-bounds}. For \(z,z'\in Y\), the corresponding
two bounds for \(B^{-1}\) and \(Q_1B^{-1}=h^{-1}q\) are
\[
 (1+2L_q/a)\sum_i2^{-i}\|z_i-z_i'\|,
 \qquad
 a^{-1}\sup_i\|q(z_i-z_i')\|,
\]
which prove the second. These are global two-point estimates, not
only estimates at the origin.

Finally, define
\[
 \iota:c_0(V)\to X,\qquad \iota((v_i))=((0,v_i))_i.
\]
The sequence belongs to \(X\), and
\[
 \|\iota((v_i))\|_X
 =\max\left\{\sum_i2^{-i}\|v_i\|,\ \sup_i\|v_i\|\right\}
 =\sup_i\|v_i\|,
\]
because \(\sum_{i\ge1}2^{-i}=1\). This is a linear isometry.
The Hilbert space \(M\) is superreflexive, so \Cref{thm:free-wsc}
makes \(E=\F(M)\) weakly sequentially complete. By hypothesis \(q\)
is completely continuous, and \Cref{lem:graph-obstruction} therefore
excludes a linear copy of \(c_0(V)\) from \(Y\). A linear isomorphism
from \(X\) onto \(Y\) would carry the displayed closed copy to one in
\(Y\), which is impossible.
\end{proof}

\begin{remark}
\label[remark]{rem:no-section}
The identities \(qB=hQ_1\) and \(Q_1B^{-1}=h^{-1}q\) involve
nonlinear maps on both sides. They do not exhibit a map \(s:U\to E\)
with \(qs=\operatorname{id}_U\), Lipschitz or otherwise. In
particular, the quotient splitting corollary of Godefroy and Kalton
\cite[Corollary~3.2]{GK03} cannot be invoked from these identities.
\end{remark}
